\chapter{Overview}
\label{chap:overview}

% archon:covers MR4688702OnTheBirationalSectionConjectureWithStrongBirationalityAssumptions.lean MR4688702OnTheBirationalSectionConjectureWithStrongBirationalityAssumptions/Basic.lean

\section{Introduction}

%%==========================================================================
%% PART 0: Notation and conventions
%%==========================================================================

% SOURCE: [Bresciani, arXiv:2108.13397], Notation and conventions, p.~4
% SOURCE QUOTE: Throughout the article, it is tacitly assumed that schemes do not have
% points of positive characteristic. In particular, fields are of characteristic $0$.
% Curves are smooth and geometrically connected. If $X$ is a curve, we denote by
% $\bar{X}$ its smooth completion. The letter $R$ always denotes a DVR with fraction
% field $K$ and residue field $k$. If $A$ is an abelian group, we denote by $\hat{A}$
% the projective limit $\projlim_{n}A/nA$. We write b.l. (resp. $t$-b.l.) as an
% abbreviation for birationally liftable (resp. $t$-birationally liftable).
\begin{remark}[Notation and conventions]
  \label{rem:notation}
  \textit{Source: Bresciani, arXiv:2108.13397, Notation and conventions.}
  Throughout, fields have characteristic \(0\). Curves are smooth and geometrically
  connected. If \(X\) is a curve, \(\bar{X}\) denotes its smooth completion. The
  letter \(R\) always denotes a DVR with fraction field \(K\) and residue field \(k\).
  If \(A\) is an abelian group, \(\hat{A}\) denotes the profinite completion
  \(\varprojlim_{n} A/nA\). We write \(\BL\) (resp.\ \(\TBL\)) as abbreviations for
  birationally liftable (resp.\ \(t\)-birationally liftable).
\end{remark}

%%==========================================================================
%% PART 1: Basic definitions
%%==========================================================================

% SOURCE: [Bresciani, arXiv:2108.13397], Introduction, p.~1
% SOURCE QUOTE: Given $X$ a geometrically connected, smooth curve over a field $k$ with
% absolute Galois group $\gal(k)$, let $\mc{S}_{X/k}$ be the set of sections of
% $\pi_{1}(X)\to\gal(k)$ modulo the action of $\pi_{1}(X_{\bar{k}})$ by conjugation,
% these are usually called Galois sections.
\begin{definition}[GaloisSection]
  \label{def:galois-section}
  \lean{BSC.TODO.GaloisSection}
  \textit{Source: Bresciani, arXiv:2108.13397, Introduction.}
  Let \(X\) be a geometrically connected smooth curve over a field \(k\) with absolute
  Galois group \(\gal(k)\). A \emph{Galois section} of \(X\) is an element of
  \(\GalSec{X}{k}\), the set of sections of \(\pi_1(X) \to \gal(k)\) modulo conjugation
  by \(\pi_1(X_{\bar{k}})\). Equivalently, in the language of the étale fundamental
  gerbe, it is an isomorphism class in \(\FundGerbe{X}{k}(k)\).
\end{definition}

% SOURCE: [Bresciani, arXiv:2108.13397], Introduction, p.~1
% SOURCE QUOTE: A Galois section is \emph{geometric} if it is associated with a
% $k$-rational point of $X$
\begin{definition}[GeometricSection]
  \label{def:geometric-section}
  \lean{BSC.TODO.GeometricSection}
  \uses{def:galois-section}
  \textit{Source: Bresciani, arXiv:2108.13397, Introduction.}
  A Galois section \(s \in \GalSec{X}{k}\) is \emph{geometric} if it is associated
  with a \(k\)-rational point of \(X\).
\end{definition}

% SOURCE: [Bresciani, arXiv:2108.13397], Introduction, p.~1
% SOURCE QUOTE: \emph{cuspidal} if it is associated with a $k$-rational point in the
% boundary $\bar{X}\setminus X$, where $\bar{X}$ is the smooth completion of $X$.
\begin{definition}[CuspidalSection]
  \label{def:cuspidal-section}
  \lean{BSC.TODO.CuspidalSection}
  \uses{def:galois-section}
  \textit{Source: Bresciani, arXiv:2108.13397, Introduction.}
  A Galois section \(s \in \GalSec{X}{k}\) is \emph{cuspidal} if it is associated
  with a \(k\)-rational point in the boundary \(\bar{X} \setminus X\), where
  \(\bar{X}\) is the smooth completion of \(X\).
\end{definition}

% SOURCE: [Bresciani, arXiv:2108.13397], Introduction, p.~2
% SOURCE QUOTE: Recall that a Galois section of $\mc{S}_{X/k}$ is \emph{birationally
% liftable} if it's in the image of $\mc{S}_{k(X)/k}\to\mc{S}_{X/k}$ \cite[\S 2.1.1]{sti15}
\begin{definition}[BiratLiftable]
  \label{def:birl}
  \lean{BSC.TODO.BiratLiftable}
  \uses{def:galois-section}
  \textit{Source: Bresciani, arXiv:2108.13397, Introduction.}
  A Galois section \(s \in \GalSec{X}{k}\) is \emph{birationally liftable} (\(\BL\))
  if it lies in the image of the natural map \(\GalSec{k(X)}{k} \to \GalSec{X}{k}\).
\end{definition}

% SOURCE: [Bresciani, arXiv:2108.13397], Introduction (Definition*), p.~3
% SOURCE QUOTE: Let $X$ be a curve over a field $k$, $s\in\mc{S}_{X/k}$ a Galois
% section and $t$ an indeterminate; consider the base change
% $s_{k(t)}\in\mc{S}_{X_{k(t)}/k(t)}$ of $s$ to $k(t)$. We say that $s$ is
% \emph{$t$-birationally liftable} if $s_{k(t)}$ is birationally liftable.
\begin{definition}[TBiratLiftable]
  \label{def:tbirl}
  \lean{BSC.TODO.TBiratLiftable}
  \uses{def:galois-section, def:birl}
  \textit{Source: Bresciani, arXiv:2108.13397, Introduction.}
  Let \(t\) be an indeterminate. A Galois section \(s \in \GalSec{X}{k}\) is
  \emph{\(t\)-birationally liftable} (\(\TBL\)) if the base change
  \(s_{k(t)} \in \GalSec{X_{k(t)}}{k(t)}\) is birationally liftable.
\end{definition}

%%==========================================================================
%% PART 2: The three conjectures
%%==========================================================================

% SOURCE: [Bresciani, arXiv:2108.13397], Introduction (SC), p.~1
% SOURCE QUOTE: For every smooth, geometrically connected hyperbolic curve $X$ over a
% field $k$ finitely generated over $\Q$, every Galois section of $X$ is either
% geometric or cuspidal.
\begin{definition}[SectionConjecture]
  \label{def:section-conjecture}
  \lean{BSC.TODO.SectionConjecturePred}
  \uses{def:galois-section, def:geometric-section, def:cuspidal-section}
  \textit{Source: Bresciani, arXiv:2108.13397, Introduction (Section Conjecture).}
  The \emph{Section Conjecture} asserts: for every smooth, geometrically connected
  hyperbolic curve \(X\) over a field \(k\) finitely generated over \(\Q\), every
  Galois section of \(X\) is either geometric or cuspidal.
\end{definition}

% SOURCE: [Bresciani, arXiv:2108.13397], Introduction (BSC), p.~2
% SOURCE QUOTE: For every smooth, geometrically connected curve $X$ over a field $k$
% finitely generated over $\Q$, every birational Galois section of $X$ is cuspidal.
\begin{definition}[BiratSectionConj]
  \label{def:bsc}
  \lean{BSC.TODO.BiratSectionConjPred}
  \uses{def:galois-section, def:cuspidal-section}
  \textit{Source: Bresciani, arXiv:2108.13397, Introduction (Birational Section Conjecture).}
  The \emph{Birational Section Conjecture} (BSC) asserts: for every smooth,
  geometrically connected curve \(X\) over a field \(k\) finitely generated over
  \(\Q\), every birational Galois section of \(X\) is cuspidal.
\end{definition}

% SOURCE: [Bresciani, arXiv:2108.13397], Introduction (CC), p.~3
% SOURCE QUOTE: For every smooth, geometrically connected hyperbolic curve $X$ over a
% field $k$ finitely generated over $\Q$ and every non-empty open subset $U\subset X$,
% the map $\mc{S}_{U/k}\to\mc{S}_{X/k}$ is surjective.
\begin{definition}[CuspidalizationConj]
  \label{def:cuspidalization}
  \lean{BSC.TODO.CuspidalizationConjPred}
  \uses{def:galois-section}
  \textit{Source: Bresciani, arXiv:2108.13397, Introduction (Cuspidalization Conjecture).}
  The \emph{Cuspidalization Conjecture} asserts: for every smooth, geometrically
  connected hyperbolic curve \(X\) over \(k\) finitely generated over \(\Q\) and
  every non-empty open \(U \subset X\), the map \(\GalSec{U}{k} \to \GalSec{X}{k}\)
  is surjective.
\end{definition}

\section{Étale fundamental gerbes}

%%==========================================================================
%% PART 3: Étale fundamental gerbes
%%==========================================================================

% NOTE: This is the fundamental gerbe of Bresciani-Vistoli (arXiv:1404.7475); absent from Mathlib.
% SOURCE: [Bresciani, arXiv:2108.13397], Section 1 (Étale fundamental gerbes), p.~5
% SOURCE QUOTE: If $X$ is geometrically connected over a field $k$, the étale fundamental
% gerbe $\Pi_{X/k}$ is a profinite étale stack over $k$ with a morphism $X\to\Pi_{X/k}$
% universal among morphisms to finite étale stacks over $k$. The space of Galois sections
% $\mc{S}_{X/k}$ of $X$ is in natural bijection with the set of isomorphism classes of
% $\Pi_{X/k}(k)$.
\begin{definition}[EtaleFundGerbe]
  \label{def:etale-fund-gerbe}
  \lean{BSC.TODO.EtaleFundGerbe}
  \uses{def:galois-section}
  \textit{Source: Bresciani, arXiv:2108.13397, Section 1.}
  % NOTE: This is the fundamental gerbe of Bresciani-Vistoli (arXiv:1404.7475); absent from Mathlib.
  For \(X\) geometrically connected over \(k\), the \emph{étale fundamental gerbe}
  \(\FundGerbe{X}{k}\) is a profinite étale stack over \(k\) equipped with a morphism
  \(X \to \FundGerbe{X}{k}\) that is universal among morphisms to finite étale stacks
  over \(k\). The set \(\GalSec{X}{k}\) of Galois sections is in natural bijection with
  the set of isomorphism classes of \(\FundGerbe{X}{k}(k)\).
\end{definition}

\subsection{Relative fundamental gerbes}

% SOURCE: [Bresciani, arXiv:2108.13397], Section 1.1 (Relative fundamental gerbes), p.~6
% SOURCE QUOTE: Let $k$ be a field of characteristic $0$ and $A$ a ring which is a
% localization of a $1$-dimensional regular domain $\tilde{A}$ of finite type over $k$,
% write $C=\spec A$. If $X\to C$ is a family of curves and $F$ is a geometric fiber,
% then $1\to\pi_{1}(F)\to\pi_{1}(X)\to\pi_{1}(C)\to 1$ is exact.
\begin{lemma}[FamiliesExactSeq]
  \label{lem:families}
  \lean{BSC.TODO.FamiliesExactSeq}
  \textit{Source: Bresciani, arXiv:2108.13397, Lemma 1 (families).}
  Let \(k\) be a field of characteristic \(0\) and \(A\) a ring which is a localization
  of a \(1\)-dimensional regular domain \(\tilde{A}\) of finite type over \(k\), write
  \(C = \spec A\). If \(X \to C\) is a family of curves and \(F\) is a geometric fiber,
  then
  \[1 \to \pi_1(F) \to \pi_1(X) \to \pi_1(C) \to 1\]
  is exact.
\end{lemma}

% SOURCE QUOTE PROOF: Assume first that $A$ is a $1$-dimensional regular domain of
% finite type over $k$. If $k=\C$, the sequence of topological fundamental groups is
% exact since the second homotopy group of $C$ is trivial, and we get an exact sequence
% of profinite completions by \cite[Proposition 5]{anderson} and
% \cite[Proposition 3.7]{gjz08}. Assume that $k$ is any field of characteristic $0$,
% it is clearly enough to do the case in which $k$ is algebraically closed. In this
% case, since everything is of finite type there exists an algebraically closed field
% $k'$ with embeddings $k'\subset k$ and $k'\subset \C$ such that $X\to C$ descends to
% a family of curves $X'\to C'$ over $k'$. Since étale fundamental groups are invariant
% under base change of algebraically closed fields \cite[Exposé X, Corollaire 1.8]{sga1},
% the statement follows from the case $k=\C$.
\begin{proof}
  Reduce to \(k = \mathbb{C}\) by invariance of étale fundamental groups under base
  change of algebraically closed fields (SGA1, Exp.\ X, Cor.\ 1.8). For \(k = \mathbb{C}\),
  the sequence of topological fundamental groups is exact since \(\pi_2(C) = 0\) (as
  \(C\) is an affine curve), and the profinite completion preserves this exactness
  (Anderson; Geraghty-Johansson-Zhu). For a localization of a finite-type domain,
  write \(\pi_1(C) \cong \varprojlim_U \pi_1(U)\) over open \(U \supset C\) in
  \(\tilde{C} = \spec \tilde{A}\), take the projective limit of the short exact
  sequences, using that the left term \(\pi_1(F)\) is constant.
\end{proof}

% SOURCE: [Bresciani, arXiv:2108.13397], Section 1.1, p.~6
% SOURCE QUOTE: Define the \emph{relative fundamental gerbe} $\Pi_{X/C}$ by the
% $2$-cartesian diagram [\Pi_{X/C} \to \Pi_{X/k}, C \to \Pi_{C/k}], there is a
% structural morphism $X\to\Pi_{X/C}$ over $C$.
\begin{definition}[RelFundGerbe]
  \label{def:rel-fund-gerbe}
  \lean{BSC.TODO.RelFundGerbe}
  \uses{def:etale-fund-gerbe, lem:families}
  \textit{Source: Bresciani, arXiv:2108.13397, Section 1.1.}
  Let \(X \to C\) be a family of curves as in \cref{lem:families}. The
  \emph{relative fundamental gerbe} \(\FundGerbe{X}{C}\) is defined by the
  \(2\)-Cartesian diagram
  \[
    \begin{array}{ccc}
      \FundGerbe{X}{C} & \to & \FundGerbe{X}{k} \\
      \downarrow & & \downarrow \\
      C & \to & \FundGerbe{C}{k},
    \end{array}
  \]
  equipped with a structural morphism \(X \to \FundGerbe{X}{C}\) over \(C\). For
  every field extension \(k'/k\) and point \(c : \spec k' \to C\), the induced
  morphism \(\FundGerbe{X_c}{k'} \to \FundGerbe{X}{C} \times_C \spec k'\) is an
  isomorphism.
\end{definition}

% SOURCE: [Bresciani, arXiv:2108.13397], Lemma 2 (limit), p.~6
% SOURCE QUOTE: The relative fundamental gerbe $\Pi_{X/C}\to C$ is a projective limit
% $\projlim \Phi_{i}$ of proper, étale morphisms $\Phi_{i}\to C$ which are gerbes over
% $C$.
\begin{lemma}[RelGerbeLimit]
  \label{lem:rel-gerbe-limit}
  \lean{BSC.TODO.RelGerbeLimit}
  \uses{def:rel-fund-gerbe}
  \textit{Source: Bresciani, arXiv:2108.13397, Lemma 2.}
  The relative fundamental gerbe \(\FundGerbe{X}{C} \to C\) is a projective limit
  \(\varprojlim \Phi_i\) of proper, étale morphisms \(\Phi_i \to C\) which are gerbes
  over \(C\).
\end{lemma}

% SOURCE QUOTE PROOF: Since $\Pi_{X/k}$, $\Pi_{C/k}$ are pro-finite étale, then we may
% write them as projective limits $\projlim_{i}\Psi_{i}$, $\projlim_{i}\Lambda_{i}$ of
% finite étale gerbes over $k$ with $\Pi_{X/k}\to\Psi_{i}$, $\Pi_{C/k}\to\Lambda_{i}$
% locally full \cite[Definition 3.4, Proposition 3.9]{bv19}. Up to re-indexing, we may
% assume that $\Pi_{X/k}\to\Pi_{C/k}$ induces a morphism $\Psi_{i}\to\Lambda_{i}$ for
% every $i$; the fact that the fibers of $X\to C$ are geometrically connected implies
% that $\Pi_{X/k}\to\Pi_{C/k}$ is locally full, which in turn implies that
% $\Psi_{i}\to\Lambda_{i}$ is locally full. The morphism $\Psi_{i}\to\Lambda_{i}$ is a
% proper étale relative gerbe. The statement then follows by defining $\Phi_{i}$ as
% $\Psi_{i}\times_{\Lambda_{i}}C$.
\begin{proof}
  \uses{def:etale-fund-gerbe}
  Write \(\FundGerbe{X}{k} = \varprojlim_i \Psi_i\) and \(\FundGerbe{C}{k} = \varprojlim_i \Lambda_i\)
  as projective limits of finite étale gerbes over \(k\) (Bresciani-Vistoli, Prop.\ 3.9).
  Geometric connectedness of the fibers of \(X \to C\) implies
  \(\FundGerbe{X}{k} \to \FundGerbe{C}{k}\) is locally full, so each
  \(\Psi_i \to \Lambda_i\) is a proper étale relative gerbe. Set
  \(\Phi_i = \Psi_i \times_{\Lambda_i} C\).
\end{proof}

\section{Non-unique specializations of ramified Galois sections}

%%==========================================================================
%% PART 4: Non-unique specializations of ramified Galois sections
%%==========================================================================

% NOTE: Construction from Talpo-Vistoli arXiv:1412.7523; absent from Mathlib.
% SOURCE: [Bresciani, arXiv:2108.13397], Section 2 (Non-unique specializations), p.~7
% SOURCE QUOTE: Let $R$ be a DVR with fraction field $K$ and residue field $k$, the
% \emph{$n$-th root stack} $\sqrt[n]{\spec R}$ of $\spec R$ is defined for every $n$.
% If $\pi\in R$ is a uniformizing parameter, then $\sqrt[n]{\spec R}$ is isomorphic to
% the quotient stack $[\spec R(\sqrt[n]{\pi})/\mu_{n}]$, but the definition of root
% stack does not depend on the choice of $\pi$. The \emph{infinite root stack}
% $\sqrt[\infty]{\spec R}$ is the projective limit $\projlim_{n}\sqrt[n]{\spec R}$.
% Let $c\in\spec R$ be the closed point and $S^{1}_{c}$ the reduced fiber of
% $\sqrt[\infty]{\spec R}$ over $c$, it is non-canonically isomorphic to the
% classifying stack $\cB_{k}\hz(1)$ of $\hz(1)$ over $k$.
\begin{definition}[InfRootStack]
  \label{def:infinite-root-stack}
  \lean{BSC.TODO.InfRootStack}
  \textit{Source: Bresciani, arXiv:2108.13397, Section 2.}
  % NOTE: Construction from Talpo-Vistoli arXiv:1412.7523; absent from Mathlib.
  For \(R\) a DVR with uniformizing parameter \(\pi\), the \emph{\(n\)-th root stack}
  \(\sqrt[n]{\spec R}\) is the quotient stack \([\spec R(\sqrt[n]{\pi})/\mu_n]\)
  (independent of the choice of \(\pi\)). The \emph{infinite root stack} is
  \(\sqrt[\infty]{\spec R} = \varprojlim_n \sqrt[n]{\spec R}\). Its closed fiber
  \(S^1_c\) (with reduced structure) is non-canonically isomorphic to the classifying
  stack \(\cB_k \hz(1)\); there is a non-canonical bijection between isomorphism
  classes of \(S^1_c(k)\) and \(\hat{k^*} = \varprojlim_n k^*/k^{*n}\).
\end{definition}

% SOURCE: [Bresciani, arXiv:2108.13397], Proposition 1 (prorootval), p.~7
% SOURCE QUOTE: Let $X\to C$ be a family of curves as in \autoref{families}, $c\in C$
% a closed point with local ring $R$ and $z:\spec K\to\Pi_{X/C}$ a generic Galois
% section, where $K$ is the fraction field of $C$. Then $z$ extends to a $2$-commutative
% diagram [\spec K \to \sqrt[\infty]{\spec R} \to \Pi_{X/C}, \spec R \to C].
% The extension $\sqrt[\infty]{\spec R}\to\Pi_{X/C}$ is unique up to a unique
% isomorphism. We call the induced morphism $S^{1}_{c}\to\Pi_{X_{c}/k(c)}$ the
% specializing loop $\gamma_{z}(c)$ of $z$ at $c$. An extension $\spec R\to\Pi_{X/C}$
% exists if and only if the specializing loop is constant.
\begin{proposition}[ProrootValCrit]
  \label{prop:prorootval}
  \lean{BSC.TODO.ProrootValCrit}
  \uses{lem:rel-gerbe-limit, def:infinite-root-stack}
  \textit{Source: Bresciani, arXiv:2108.13397, Proposition 1.}
  Let \(X \to C\) be a family of curves as in \cref{lem:families}, \(c \in C\) a
  closed point with local ring \(R\), and \(z : \spec K \to \FundGerbe{X}{C}\) a
  generic Galois section (\(K\) = fraction field of \(C\)). Then \(z\) extends
  uniquely (up to unique isomorphism) to a morphism \(\sqrt[\infty]{\spec R} \to \FundGerbe{X}{C}\)
  compatible with \(\spec K \hookrightarrow \sqrt[\infty]{\spec R}\) and the map
  \(\spec R \to C\). An extension \(\spec R \to \FundGerbe{X}{C}\) exists if and only
  if the induced specializing loop \(S^1_c \to \FundGerbe{X_c}{k(c)}\) is constant.
\end{proposition}

% SOURCE QUOTE PROOF: Using \autoref{limit}, this is a direct consequence of
% \cite[Theorem 3.1]{giulio-angelo-valuative}.
\begin{proof}
  \uses{lem:rel-gerbe-limit, def:infinite-root-stack}
  By \cref{lem:rel-gerbe-limit}, \(\FundGerbe{X}{C} \to C\) is a projective limit of
  proper morphisms; the claim then follows from the valuative criterion for pro-finite
  proper stacks (Bresciani-Vistoli, arXiv:1404.7475, Theorem 3.1).
\end{proof}

% SOURCE: [Bresciani, arXiv:2108.13397], Definition 1 (loopdef), p.~8
% SOURCE QUOTE: A \emph{specialization} of $z$ at $c$ is any Galois section of
% $X_{c}/k(c)$ in the essential image of the specializing loop
% $S^{1}_{c}(k(c))\to\Pi_{X_{c}/k(c)}(k(c))$ of $z$ at $c$ evaluated on $k(c)$.
% A specialization always exists, but in general it might be not unique.
\begin{definition}[SpecLoop]
  \label{def:spec-loop}
  \lean{BSC.TODO.SpecLoop}
  \uses{prop:prorootval}
  \textit{Source: Bresciani, arXiv:2108.13397, Definition 1.}
  With notation as in \cref{prop:prorootval}, the \emph{specializing loop} of \(z\)
  at \(c\) is the induced morphism \(\gamma_z(c) : S^1_c \to \FundGerbe{X_c}{k(c)}\).
  A \emph{specialization} of \(z\) at \(c\) is any Galois section of \(X_c/k(c)\) in
  the essential image of \(\gamma_z(c)(k(c)) : S^1_c(k(c)) \to \FundGerbe{X_c}{k(c)}(k(c))\).
  A specialization always exists but need not be unique; it is unique precisely when
  \(\gamma_z(c)\) is constant.
\end{definition}

% SOURCE: [Bresciani, arXiv:2108.13397], Example 1 (gmloop), p.~8
% SOURCE QUOTE: Let $R$ be a DVR as in \autoref{families} with fraction field $K$ and
% valuation $v:K^{*}\to\Z$, this extends naturally to an homomorphism
% $\hat{K^{*}}\to\hz$ which we still call $v$. The morphism
% $\G_{m}=\A^{1}_{R}\setminus\{0\}\to \spec R$ is a family of curves. The section
% $1\in\G_{m}(R)$ gives an identification $\Pi_{\G_{m}/R}=\cB_{R}\hz(1)$, in
% particular $\Pi_{\G_{m}/R}(R)=\hat{R^*}$ and $\Pi_{\G_{m,K}/K}(K)=\hat{K^*}$.
% If $z\in\hat{K^{*}}$ is a generic section, by \autoref{prorootval} the specializing
% loop $\gamma_{z}$ is constant if and only if $z$ is in the image of
% $\hat{R^{*}}\to\hat{K^{*}}$, or equivalently if and only if $v(z)=0\in\hz$.
\begin{example}[GmLoop]
  \label{example:gm-loop}
  \lean{BSC.TODO.GmLoop}
  \uses{def:spec-loop, prop:prorootval}
  \textit{Source: Bresciani, arXiv:2108.13397, Example 1.}
  Let \(R\) be a DVR as in \cref{lem:families} with fraction field \(K\) and
  valuation \(v : K^* \to \Z\) (extended to \(\hat{K^*} \to \hz\)). The family
  \(\Gm = \A^1_R \setminus \{0\} \to \spec R\) has
  \(\FundGerbe{\Gm}{R} \cong \cB_R \hz(1)\), so
  \(\FundGerbe{\Gm_K}{K}(K) = \hat{K^*}\). For \(z \in \hat{K^*}\) a generic
  section, the specializing loop \(\gamma_z\) is constant if and only if
  \(z \in \hat{R^*}\), i.e.\ \(v(z) = 0 \in \hz\).
\end{example}

% SOURCE: [Bresciani, arXiv:2108.13397], Lemma 3 (uniqueloop), p.~8
% SOURCE QUOTE: Let $T$ be a torus over a field $k$ with a surjective valuation
% $\nu:k^{*}\to\Z$. Suppose that we have a morphism of gerbes $\gamma:\cB_{k}\hz(1)\to
% \Pi_{T/k}$. If the image of $\cB_{k}\hz(1)(k)$ in $\Pi_{T/k}(k_{\nu})$ has a finite
% number of isomorphism classes, then $\gamma$ is constant.
\begin{lemma}[UniqueLoopTorus]
  \label{lem:unique-loop-torus}
  \lean{BSC.TODO.UniqueLoopTorus}
  \uses{def:etale-fund-gerbe}
  \textit{Source: Bresciani, arXiv:2108.13397, Lemma 3.}
  Let \(T\) be a torus over \(k\) with surjective valuation \(\nu : k^* \to \Z\),
  and let \(\gamma : \cB_k \hz(1) \to \FundGerbe{T}{k}\) be a morphism of gerbes.
  If the image of \(\cB_k \hz(1)(k)\) in \(\FundGerbe{T}{k}(k_\nu)\) has finitely
  many isomorphism classes, then \(\gamma\) is constant.
\end{lemma}

% SOURCE QUOTE PROOF: Using the image of the preferred $k$-rational section of
% $\cB_{k}\hz(1)$, we may identify $\Pi_{T/k}$ with $\mc{B}_{k}\upi_{1}(T)$, and we
% reduce to prove the following: if $f:\hz(1)\to\pi_{1}(T_{\bar{k}})$ is a Galois
% equivariant homomorphism such that the composition
% $\H^{1}(k,\hz(1))\xar{f} \H^{1}(k,\pi_{1}(T_{\bar{k}})) \xar{res}
% \H^{1}(k_{\nu},\pi_{1}(T_{\bar{k}}))$ has finite image, then $f$ is trivial.
% Choose $k'/k$ a finite extension such that $T_{k'}\simeq\G_{m}^{n}$. The valuation
% $\nu$ induces $\H^{1}(k_{\nu},\hz(1)) \to \hz$. The hypothesis implies $f \circ r = rf = 0$;
% since $r \neq 0$ and $\hz$ has no torsion, $f = 0$.
\begin{proof}
  Choose a finite extension \(k'/k\) splitting \(T\), so \(T_{k'} \cong \Gm^n\) and
  \(\pi_1(T_{\bar{k}}) \cong \hz^n(1)\). The surjective valuation \(\nu\) induces a
  surjection \(H^1(k, \hz(1)) \to \hz\). Finiteness of the image forces the composite
  \(H^1(k, \hz(1)) \xrightarrow{f} H^1(k_\nu, \pi_1(T_{\bar{k}}))\) to have finite
  image, hence the map \(\hz \to \hz^n\) given by \(rf\) is zero. Since \(r \neq 0\)
  and \(\hz\) is torsion-free, \(f = 0\).
\end{proof}

% SOURCE: [Bresciani, arXiv:2108.13397], Lemma 4 (uniqueloop1), p.~9
% SOURCE QUOTE: Let $k$ be a number field with a finite place $\nu$, $X$ a curve over
% $k$, $\gamma:\cB_{k}\hz(1)\to\Pi_{X/k}$ a morphism of gerbes. If the image of
% $\cB_{k}\hz(1)(k)$ in $\Pi_{X/k}(k_{\nu})$ has a finite number of isomorphism
% classes, then $\gamma$ is constant.
\begin{lemma}[UniqueLoopCurve]
  \label{lem:unique-loop-curve}
  \lean{BSC.TODO.UniqueLoopCurve}
  \uses{lem:unique-loop-torus}
  \textit{Source: Bresciani, arXiv:2108.13397, Lemma 4.}
  Let \(k\) be a number field, \(\nu\) a finite place, \(X\) a curve over \(k\), and
  \(\gamma : \cB_k \hz(1) \to \FundGerbe{X}{k}\) a morphism of gerbes. If the image
  of \(\cB_k \hz(1)(k)\) in \(\FundGerbe{X}{k}(k_\nu)\) has finitely many isomorphism
  classes, then \(\gamma\) is constant.
\end{lemma}

% SOURCE QUOTE PROOF: If by contradiction $\gamma$ is not constant, up to replacing
% $X$ with a finite étale cover we may assume that the composition
% $\cB_{k}\hz(1)\to\Pi_{X/k}\to\Pi_{X/k}^{\rm ab}$ is not constant. Let $J$ be the
% semi-abelian Jacobian of $X$, it is an extension of an abelian variety $A$ by a
% torus $T$. By weight reasons, the composition $\hz(1)\to\on{T}J\to\on{T}A$ is
% trivial, hence we have a factorization $\hz(1)\to\on{T}T\to\on{T}J$. By Mattuck's
% theorem \cite{mat55} $A(k_{\nu})$ has finite torsion, thus
% $\H^{1}(k_{\nu},\on{T}T)\to\H^{1}(k_{\nu},\on{T}J)$ is injective. The hypothesis
% implies that the image of $\H^{1}(k,\hz(1))$ in $\H^{1}(k_{\nu},\on{T}T)$ is
% finite, hence $\hz(1)\to\on{T}T$ is trivial by \autoref{uniqueloop}, contradiction.
\begin{proof}
  \uses{lem:unique-loop-torus}
  If \(\gamma\) is non-constant, pass to a finite étale cover and assume the
  abelianization \(\cB_k \hz(1) \to \FundGerbe{X}{k}^{\mathrm{ab}} = \FundGerbe{J}{k}\)
  is non-constant, where \(J\) is the semi-abelian Jacobian. By weight reasons
  \(\hz(1) \to TJ\) factors through the torus part \(TT\). By Mattuck's theorem
  \(A(k_\nu)\) has finite torsion, so \(H^1(k_\nu, TT) \to H^1(k_\nu, TJ)\) is
  injective. The finiteness hypothesis then forces \(\hz(1) \to TT\) to be trivial
  by \cref{lem:unique-loop-torus}, a contradiction.
\end{proof}

\section{Properties of \(t\)-birationally liftable Galois sections}

%%==========================================================================
%% PART 5: Properties of t-birationally liftable sections
%%==========================================================================

% SOURCE: [Bresciani, arXiv:2108.13397], Section 3 (Properties of t-b.l. sections),
% first lemma (unnamed), p.~10
% SOURCE QUOTE: Let $R$ be a DVR, $X\to \spec R$ a family of curves and $p\in X$ a
% closed point in the closed fiber. There exists a non-empty divisor $D\subset X$
% finite étale over $R$ with $p\in D$.
\begin{lemma}[FinEtaleDivisor]
  \label{lem:fin-etale-divisor}
  \lean{BSC.TODO.FinEtaleDivisor}
  \textit{Source: Bresciani, arXiv:2108.13397, Section 3, first lemma.}
  Let \(R\) be a DVR, \(X \to \spec R\) a family of curves, and \(p \in X\) a closed
  point in the closed fiber. There exists a non-empty divisor \(D \subset X\) finite
  étale over \(R\) with \(p \in D\).
\end{lemma}

% SOURCE QUOTE PROOF: Let $K,k$ be the fraction and residue fields of $R$. If $p$ is
% $k$-rational, this essentially follows from Bertini's theorem by taking a generic
% hyperplane section passing through $p$. If $p$ is not $k$-rational, we use a
% Bertini-like argument to reduce to the case of $\P^{1}_{R}$, where we are able to
% make an explicit construction. [...] Since we are in characteristic $0$, then
% $k(p)\simeq k[x]/\bar{q}(x)$ for some irreducible, separable polynomial
% $\bar{q}\in k[x]$, let $q\in R[x]$ be any lifting of $\bar{q}$. By construction,
% $R'=R[x]/q(x)$ is finite étale over $R$ and has only one closed point, hence it is
% a DVR. [...] there are infinitely many points of $\P^{1}(K')$ which specialize to
% $p\in\P^{1}(k(p))$, hence we may choose one in $U(K')$; denote by $E\subset \P^{1}_{R}$
% its closure [...] $E$ is étale over $R$.
\begin{proof}
  Embed \(\bar{X} \subset \PP^n_R\) and choose a generic \((n-2)\)-dimensional linear
  subspace \(L_k \subset \PP^n_k\) not meeting \(\bar{X}_k\), such that the
  hyperplane \(H_i\) spanned by \(L_k\) and each geometric point \(p_i\) over \(p\)
  meets \(\bar{X}_{\bar{k}}\) transversally and avoids \(D_0\). The projection from
  \(L\) gives \(f : \bar{X} \to \PP^1_R\) étale near \(f(p)\). Since
  \(k(p) \cong k[x]/\bar{q}(x)\) for a separable \(\bar{q}\), lift to
  \(R' = R[x]/q(x)\) (a DVR, finite étale over \(R\)); choose a point
  \(E \subset \PP^1_R\) finite étale over \(R\) through \(f(p)\); then
  \(f^{-1}(E)\) is the desired divisor.
\end{proof}

% SOURCE: [Bresciani, arXiv:2108.13397], Corollary 1 (cut), p.~13
% SOURCE QUOTE: Let $R$ be a DVR, $X\to \spec R$ a family of curves. There exists a
% direct system $(D_{i})_{i}$ of divisors finite étale over $R$ such that
% $X_{k}\setminus\bigcup_{i}D_{i}$ contains only the generic point of $X_{k}$.
\begin{corollary}[CutLemma]
  \label{cor:cut}
  \lean{BSC.TODO.CutLemma}
  \uses{lem:fin-etale-divisor}
  \textit{Source: Bresciani, arXiv:2108.13397, Corollary 1.}
  Let \(R\) be a DVR and \(X \to \spec R\) a family of curves. There exists a direct
  system \((D_i)_i\) of divisors finite étale over \(R\) such that
  \(X_k \setminus \bigcup_i D_i\) contains only the generic point of \(X_k\).
\end{corollary}

\begin{proof}
  \uses{lem:fin-etale-divisor}
  Apply \cref{lem:fin-etale-divisor} to each closed point of \(X_k\); the resulting
  direct system of divisors exhausts all closed points.
\end{proof}

% SOURCE: [Bresciani, arXiv:2108.13397], Lemma 5 (spebir), p.~13
% SOURCE QUOTE: Specializations of b.l. sections are b.l.
\begin{lemma}[SpecBirat]
  \label{lem:spec-birl}
  \lean{BSC.TODO.SpecBirat}
  \uses{cor:cut, prop:prorootval, def:birl}
  \textit{Source: Bresciani, arXiv:2108.13397, Lemma 5.}
  Specializations of \(\BL\) sections are \(\BL\). More precisely: let
  \(X \to \spec R\) be a family of curves and \(z \in \FundGerbe{X_K}{K}(K)\) a
  generic \(\BL\) Galois section. Then every specialization of \(z\) is \(\BL\).
\end{lemma}

% SOURCE QUOTE PROOF: Let $D_{i}\s X$ be a direct system of divisors as in
% \autoref{cut}. We have that $X_{i}=X\setminus D_{i}$ is a family of curves for
% every $i$, write $X_{\infty}=\projlim_{i}X_{i}$ and
% $\Pi_{X_{\infty}/R}=\projlim_{i}\Pi_{X_{i}/R}$. The closed fiber of $X_{\infty}$
% is naturally isomorphic to $\spec k(X_{k})$. Since $z$ is b.l. and
% $\spec K(X)\s X_{\infty,K}$, there exists a lifting $z'\in\Pi_{X_{\infty,K}/K}(K)$
% of $z$ to $X_{\infty,K}$. By a limit argument, the specializing loops of $z'$ in
% $\Pi_{X_{i}/R}$ induce a specializing loop $\gamma_{z'}:S^{1}_{c}\to
% \Pi_{X_{\infty,k}/k}=\Pi_{k(X_{k})/k}$, and the composition of $\gamma_{z'}$ with
% $\Pi_{k(X_{k})/k}\to \Pi_{X_{k}/k}$ is isomorphic to $\gamma_{z}$. Hence, every
% specialization of $z$ lifts to $\spec k(X_{k})$.
\begin{proof}
  \uses{cor:cut, prop:prorootval}
  Let \((D_i)\) be as in \cref{cor:cut}; set \(X_i = X \setminus D_i\) and
  \(X_\infty = \varprojlim_i X_i\), whose closed fiber is \(\spec k(X_k)\).
  Since \(z\) is \(\BL\), it lifts to \(z' \in \FundGerbe{X_{\infty,K}}{K}(K)\).
  By a limit argument with \cref{prop:prorootval}, the specializing loops of \(z'\)
  induce \(\gamma_{z'} : S^1_c \to \FundGerbe{k(X_k)}{k}\). Every specialization of
  \(z\) thus lifts to \(\spec k(X_k)\), so it is \(\BL\).
\end{proof}

% SOURCE: [Bresciani, arXiv:2108.13397], Corollary 2 (tbb), p.~13
% SOURCE QUOTE: $t$-b.l. sections are b.l.
\begin{corollary}[TBiratBirat]
  \label{cor:tbirl-birl}
  \lean{BSC.TODO.TBiratBirat}
  \uses{lem:spec-birl, def:tbirl}
  \textit{Source: Bresciani, arXiv:2108.13397, Corollary 2.}
  Every \(\TBL\) Galois section is \(\BL\).
\end{corollary}

% SOURCE QUOTE PROOF: Let $X$ be a curve over a field $k$ and $s\in\Pi_{X/k}(k)$ a
% $t$-b.l. Galois section. Write $R=k[t]_{(t)}$, we have that $s_{k(t)}$ is a generic
% section of $\Pi_{X_{R}/R}$ and $s$ is a specialization of $s_{k(t)}$. By hypothesis
% $s_{k(t)}$ is b.l., hence its specialization $s$ is b.l. by \autoref{spebir}.
\begin{proof}
  \uses{lem:spec-birl}
  Write \(R = k[t]_{(t)}\); then \(s_{k(t)}\) is a generic section of
  \(\FundGerbe{X_R}{R}\) and \(s\) is a specialization of \(s_{k(t)}\). Since
  \(s_{k(t)}\) is \(\BL\) by hypothesis, \cref{lem:spec-birl} gives that \(s\) is
  \(\BL\).
\end{proof}

% SOURCE: [Bresciani, arXiv:2108.13397], Lemma 6 (spetbir), p.~14
% SOURCE QUOTE: Specializations of $t$-b.l. sections are $t$-b.l.
\begin{lemma}[SpecTBirat]
  \label{lem:spec-tbirl}
  \lean{BSC.TODO.SpecTBirat}
  \uses{lem:spec-birl, def:tbirl}
  \textit{Source: Bresciani, arXiv:2108.13397, Lemma 6.}
  Specializations of \(\TBL\) sections are \(\TBL\). Precisely: let
  \(X \to \spec R\) be a family of curves, \(z\) a generic \(\TBL\) section, and
  \(s\) a specialization of \(z\). Then \(s\) is \(\TBL\).
\end{lemma}

% SOURCE QUOTE PROOF: Let $X\to \spec R$ be as above, $z\in\Pi_{X_{K}/K}(K)$ a
% generic, $t$-b.l. Galois section, $s\in\Pi_{X_{k}/k}(k)$ a specialization.
% Let $R'$ be the local ring of the generic point of the divisor
% $\A^{1}_{k}\s\A^{1}_{R}$. We have that $R'$ is a DVR with fraction field $K'=K(t)$
% and residue field equal $k'=k(t)$. Consider the family of curves $X_{R'}\to \spec R'$,
% then $z_{K(t)}=z_{K'}\in \Pi_{X_{K'}/K'}(K')$ is by hypothesis a b.l. Galois section.
% We have that $s_{k(t)}=s_{k'}$ is a specialization of $z_{K'}$ and thus it is b.l.
% by \autoref{spebir}.
\begin{proof}
  \uses{lem:spec-birl}
  Let \(R'\) be the local ring of the generic point of \(\A^1_k \subset \A^1_R\); it
  is a DVR with fraction field \(K(t)\) and residue field \(k(t)\). In the family
  \(X_{R'} \to \spec R'\), the base change \(z_{K(t)}\) is \(\BL\) by hypothesis.
  Then \(s_{k(t)}\) is a specialization of \(z_{K(t)}\), hence \(\BL\) by
  \cref{lem:spec-birl}, so \(s\) is \(\TBL\).
\end{proof}

% SOURCE: [Bresciani, arXiv:2108.13397], Lemma 7 (morb), p.~14
% SOURCE QUOTE: Let $f:Y\to X$ be a dominant morphism of curves over a field $k$ and
% $s\in\Pi_{Y/k}(k)$ a Galois section. If $s$ is b.l. then $f(s)\in\Pi_{X/k}(k)$ is
% b.l. If $f$ is finite étale, the converse holds.
\begin{lemma}[MorphBirat]
  \label{lem:morph-birl}
  \lean{BSC.TODO.MorphBirat}
  \uses{def:birl}
  \textit{Source: Bresciani, arXiv:2108.13397, Lemma 7.}
  Let \(f : Y \to X\) be a dominant morphism of curves over \(k\) and
  \(s \in \FundGerbe{Y}{k}(k)\) a Galois section. If \(s\) is \(\BL\), then
  \(f(s) \in \FundGerbe{X}{k}(k)\) is \(\BL\). If \(f\) is finite étale, the
  converse holds.
\end{lemma}

% SOURCE QUOTE PROOF: The first statement is obvious. If $f$ is finite étale, the
% second statement follows from the fact that
% $\gal(k(Y))=\pi_{1}(Y)\times_{\pi_{1}(X)}\gal(k(X))$.
\begin{proof}
  The forward direction is immediate. For the converse when \(f\) is finite étale,
  use the fibre product identity
  \(\gal(k(Y)) = \pi_1(Y) \times_{\pi_1(X)} \gal(k(X))\).
\end{proof}

% SOURCE: [Bresciani, arXiv:2108.13397], Corollary 3 (mortb), p.~14
% SOURCE QUOTE: Let $f:Y\to X$ be a dominant morphism of curves over a field $k$ and
% $s\in\Pi_{Y/k}(k)$ a Galois section. If $s$ is $t$-b.l. then $f(s)\in\Pi_{X/k}(k)$
% is $t$-b.l. If $f$ is finite étale, the converse holds.
\begin{corollary}[MorphTBirat]
  \label{cor:morph-tbirl}
  \lean{BSC.TODO.MorphTBirat}
  \uses{lem:morph-birl, def:tbirl}
  \textit{Source: Bresciani, arXiv:2108.13397, Corollary 3.}
  Same as \cref{lem:morph-birl} with \(\BL\) replaced by \(\TBL\): if \(f\) is
  dominant and \(s\) is \(\TBL\), then \(f(s)\) is \(\TBL\); if \(f\) is finite
  étale, the converse holds.
\end{corollary}

% SOURCE QUOTE PROOF: Follows directly from \autoref{morb} by base change to $k(t)$.
\begin{proof}
  \uses{lem:morph-birl}
  Base-change to \(k(t)\): \(s\) is \(\TBL\) iff \(s_{k(t)}\) is \(\BL\).
  Applying \cref{lem:morph-birl} to \(f_{k(t)}\) and \(s_{k(t)}\) gives
  \(f(s_{k(t)}) = f(s)_{k(t)}\) is \(\BL\), hence \(f(s)\) is \(\TBL\).
  If \(f\) is finite étale, the converse follows by the same base-change argument
  using the converse direction of \cref{lem:morph-birl}.
\end{proof}

% SOURCE: [Bresciani, arXiv:2108.13397], Lemma 8 (nflift), p.~15
% SOURCE QUOTE: Let $X$ be a curve over a number field $k$, $U\s X$ a non-empty open
% subset and $s\in\Pi_{X/k}(k)$ a $t$-b.l. Galois section. There exists a $t$-b.l.
% section of $U$ which lifts $s$.
\begin{lemma}[NFLift]
  \label{lem:nf-lift}
  \lean{BSC.TODO.NFLift}
  \uses{cor:tbirl-birl, lem:unique-loop-curve, prop:prorootval}
  \textit{Source: Bresciani, arXiv:2108.13397, Lemma 8.}
  Let \(X\) be a curve over a number field \(k\), \(U \subseteq X\) a non-empty open
  subset, and \(s \in \FundGerbe{X}{k}(k)\) a \(\TBL\) Galois section. There exists
  a \(\TBL\) section of \(U\) which lifts \(s\).
\end{lemma}

% SOURCE QUOTE PROOF: We may assume that $X$ is non-parabolic and that $s$ is neither
% geometric nor cuspidal, otherwise this is trivial. By hypothesis, $s_{k(t)}$ lifts to
% a b.l. Galois section $r_{0}$ of $U_{k(t)}$, we are going to show that $r_{0}$ is
% the base change to $k(t)$ of some Galois section of $U$. [...] For every finite place
% $\nu$ of $k$ there is a unique local point $x_{\nu}\in\bar{X}(k_{\nu})$ associated
% with $s$. Since $s$ lifts to a b.l. section of $U$ by \autoref{tbb} and it is not
% geometric nor cuspidal, the set of places $\nu$ such that $x_{\nu}\in U$ has Dirichlet
% density $1$ by \cite[Theorem B]{sti15}, hence we may choose $\nu$ with $x_{\nu}\in U$
% and $k'\subset k_{\nu}$. If $r_{0,c}$ is a specialization of $r_{0}$ at $c$, by
% construction it is b.l. and lifts $s$, hence $r_{0,c,k_{\nu}}$ is geometric associated
% with $x_{\nu}$. By \autoref{uniqueloop1}, the specializing loop
% $S^{1}_{c}\to\Pi_{U_{k'}/k'}$ of $r_{0}$ at $c$ is constant and hence $r_{0}$ extends
% to a section $\spec\O_{c}\to\Pi_{U\times\A^{1}/\A^{1}}$ by \autoref{prorootval}. [...]
% by \cite[Corollary A.3]{fed} we get that $r_{0}$ extends to a section
% $\tilde{r}_{0}:\A^{1}\to\Pi_{U\times\A^{1}/\A^{1}}$.
\begin{proof}
  \uses{cor:tbirl-birl, lem:unique-loop-curve, prop:prorootval, lem:spec-birl}
  Assume \(X\) non-parabolic and \(s\) not geometric/cuspidal. By hypothesis,
  \(s_{k(t)}\) lifts to a \(\BL\) section \(r_0\) of \(U_{k(t)}\). For each closed
  point \(c \in \A^1\) with residue field \(k'\), choose a finite place \(\nu\) of
  \(k\) with \(k' \subset k_\nu\) and \(x_\nu \in U\) (possible by Stix's Theorem~B,
  which gives density~1). Any specialization \(r_{0,c}\) is \(\BL\) by
  \cref{lem:spec-birl} and maps to \(s\); over \(k_\nu\) it is the geometric section
  at \(x_\nu \in U \setminus \{c\}\). By \cref{lem:unique-loop-curve}, the
  specializing loop of \(r_0\) at \(c\) is constant, so \(r_0\) extends over
  \(\mathcal{O}_c\) by \cref{prop:prorootval}. Doing this for every closed point and
  using that \(\FundGerbe{U \times \A^1}{\A^1} \cong \FundGerbe{U}{k} \times_k \A^1\)
  (Bresciani, arXiv:2003.04649, Cor.\ A.3), we get \(r_0 = r_{k(t)}\) for a
  \(\TBL\) section \(r\) of \(U\) lifting \(s\).
\end{proof}

\section{The main argument}

%%==========================================================================
%% PART 6: The main argument
%%==========================================================================

% SOURCE: [Bresciani, arXiv:2108.13397], Proposition 2 (simple), p.~16
% SOURCE QUOTE: Let $k$ be a number field, $\nu$ a finite place, write
% $U=\A^{1}\setminus\{1,2\}$. Let $s\in\Pi_{U/k}(k)$ be a b.l. Galois section and
% $x_{\nu}\in\P^{1}(k_{\nu})$ the unique associated local point. Let
% $\delta:\spec k(t)\to\A^{1}$ be the ``diagonal'' point, i.e. the generic one, and
% assume that $s_{k(t)}$ lifts to a section of
% $\spec k(\A^{1})\otimes_{k} k(t)\setminus\{\delta\}$ (e.g. if $s$ is $t$-b.l.).
% Then $x_{\nu}$ is $k$-rational.
\begin{proposition}[SimpleProp]
  \label{prop:simple}
  \lean{BSC.TODO.SimpleProp}
  \uses{lem:unique-loop-curve, example:gm-loop, prop:prorootval, def:birl}
  \textit{Source: Bresciani, arXiv:2108.13397, Proposition 2.}
  Let \(k\) be a number field, \(\nu\) a finite place, and \(U = \A^1 \setminus \{1,2\}\).
  Let \(s \in \FundGerbe{U}{k}(k)\) be a \(\BL\) Galois section and
  \(x_\nu \in \PP^1(k_\nu)\) the unique associated local point. Let
  \(\delta : \spec k(t) \to \A^1\) be the diagonal (generic) point, and assume that
  \(s_{k(t)}\) lifts to a section of
  \(\spec k(\A^1) \otimes_k k(t) \setminus \{\delta\}\) (e.g.\ if \(s\) is \(\TBL\)).
  Then \(x_\nu\) is \(k\)-rational.
\end{proposition}

% SOURCE QUOTE PROOF: We may assume that $x_{\nu}\in U$ [...] We divide the proof in
% three steps. Step 1. Specializations of $r$. [...] the specializing loop of $r$ is
% constant: [...] $r_{c,k_{\nu}}$ is \emph{the} geometric section associated to
% $x_{\nu}$ [...] By \autoref{uniqueloop1}, the specializing loop of $r$ at $c$ is
% constant. Step 2. Change of coordinates. The map $y\mapsto t-y$ defines an
% automorphism $\phi:\A^{1}_{k(t)}\to\A^{1}_{k(t)}$ with $\phi(\delta)=0$ [...] we
% get a Galois section $\phi(s_{k(t)})\in\Pi_{V}(k(t))$ with a lifting
% $\phi(r)\in\Pi_{V\setminus \{0\}}(k(t))$ [...] $z=\lambda\cdot q^{e}\cdot(t-1)^{e_{1}}
% \cdot(t-2)^{e_{2}}$ [...] Step 3. Specializations of $z$. [...] Using a sequence
% $c_{i}\to x_{\eta}$ [...] $e=1$ [...] if by contradiction $x_{\nu}$ is not
% $k$-rational and thus $\deg q\ge 2$, using a sequence $c_{i}\to y_{j}$ we see that
% $\eta(\lambda)+\sum_{j}\eta(c-y_{j})=0\in\hz$ is not constant, which is absurd.
\begin{proof}
  \uses{lem:unique-loop-curve, example:gm-loop, prop:prorootval, def:birl}
  \textbf{Step 1 (Specializations of \(r\)).}
  Let \(r_0\) be the lifting and \(r\) its image in \(U_{k(t)} \setminus \{\delta\}\).
  For each \(k\)-rational \(c \in U(k)\) with \(c \neq x_\nu\), any specialization
  \(r_c\) of \(r\) at \(c\) is \(\BL\) and maps to \(s\); over \(k_\nu\) it must be
  the geometric section at \(x_\nu \in U \setminus \{c\}\). By
  \cref{lem:unique-loop-curve}, the specializing loop at \(c\) is constant.

  \textbf{Step 2 (Change of coordinates).}
  The automorphism \(\phi : y \mapsto t - y\) sends \(U_{k(t)} \xrightarrow{\sim} V\)
  where \(V = \A^1_{k(t)} \setminus \{t-1, t-2\}\). Then
  \(\phi(r) \in \FundGerbe{V \setminus \{0\}}{}(k(t)) \hookrightarrow \hat{k(t)^*}\)
  gives \(z = \lambda \cdot q^e \cdot (t-1)^{e_1} \cdot (t-2)^{e_2}\) where \(q\) is
  the minimal polynomial of \(x_\nu\), \(e, e_1, e_2 \in \hz\), \(\lambda \in \hat{k^*}\).
  Step 1 and \cref{example:gm-loop} force \(e_c = 0\) for all closed \(c \neq x_\nu, 1, 2\).

  \textbf{Step 3 (Specializations of \(z\)).}
  Fix a finite extension \(K\) splitting \(q\) and extension \(\eta\) of \(\nu\).
  Taking sequences \(c_i \to 1\) and \(c_i \to 2\) forces \(e_1 = e_2 = 0\).
  Taking \(c_i \to x_\eta\) forces \(e = 1\). If \(\deg q \geq 2\), taking
  \(c_i \to y_j\) (a conjugate of \(x_\nu\)) gives a contradiction, so
  \(\deg q = 1\) and \(x_\nu \in k\).
\end{proof}

\section{Reduction to number fields}

%%==========================================================================
%% PART 7: Reduction to number fields
%%==========================================================================

% SOURCE: [Bresciani, arXiv:2108.13397], Section 4 (Reduction to number fields), p.~18
% SOURCE QUOTE: Let $X$ be a geometrically connected curve over a field $k$. A Galois
% section $s$ of $X$ is quasi-$t$-b.l. if there is an open subset $U\s X$, a lifting
% $s'$ of $s$ to $U$ and a dominant map $f:U\to Y$ where $Y$ is a non-parabolic curve
% such that $s'$ is b.l. and $f(s')$ is $t$-b.l.
\begin{definition}[QuasiTBirat]
  \label{def:quasi-tbirl}
  \lean{BSC.TODO.QuasiTBirat}
  \uses{def:tbirl, def:birl}
  \textit{Source: Bresciani, arXiv:2108.13397, Section 4.}
  Let \(X\) be a geometrically connected curve over \(k\). A Galois section \(s\) of
  \(X\) is \emph{quasi-\(\TBL\)} if there exists an open subset \(U \subseteq X\),
  a lifting \(s'\) of \(s\) to \(U\), and a dominant map \(f : U \to Y\) with \(Y\)
  a non-parabolic curve, such that \(s'\) is \(\BL\) and \(f(s')\) is \(\TBL\).
\end{definition}

% SOURCE: [Bresciani, arXiv:2108.13397], Lemma 9 (spequasi), p.~18
% SOURCE QUOTE: Let $C$ be a curve with fraction field $K$, $X\to C$ a family of curves
% and $s\in\Pi_{X_{K}/K}(K)$ a generic section which is quasi-$t$-b.l. There exists a
% non-empty open subset $U\s C$ such that the specializations of $X|_{U}\to U$ are
% quasi-$t$-b.l.
\begin{lemma}[SpecQuasi]
  \label{lem:spec-quasi}
  \lean{BSC.TODO.SpecQuasi}
  \uses{def:quasi-tbirl, lem:spec-birl, lem:spec-tbirl}
  \textit{Source: Bresciani, arXiv:2108.13397, Lemma 9.}
  Let \(C\) be a curve with fraction field \(K\), \(X \to C\) a family of curves, and
  \(s \in \FundGerbe{X_K}{K}(K)\) a generic quasi-\(\TBL\) section. There exists a
  non-empty open \(U \subseteq C\) such that the specializations of \(X|_U \to U\) are
  quasi-\(\TBL\).
\end{lemma}

% SOURCE QUOTE PROOF: Up to shrinking both $C$ and $X$, we may assume that there exists
% a family of non-parabolic curves $Y\to C$ and a quasi-finite morphism $f:X\to Y$ over
% $C$ such that $f(s)$ is $t$-b.l. The specializations of $f(s)$ are $t$-b.l. by
% \autoref{spetbir} and the specializations of $s$ are b.l. by \autoref{spebir}, it
% follows that they are quasi-$t$-b.l., too.
\begin{proof}
  \uses{lem:spec-birl, lem:spec-tbirl}
  Shrink \(C\) and \(X\) to arrange a non-parabolic family \(Y \to C\) and
  quasi-finite \(f : X \to Y\) with \(f(s)\) \(\TBL\). Specializations of \(f(s)\)
  are \(\TBL\) by \cref{lem:spec-tbirl}; specializations of \(s\) are \(\BL\) by
  \cref{lem:spec-birl}; hence they are quasi-\(\TBL\).
\end{proof}

% SOURCE: [Bresciani, arXiv:2108.13397], Lemma 10 (quasinf), p.~18
% SOURCE QUOTE: Let $k$ be a number field, and assume that $t$-b.l. sections over $k$
% are geometric or cuspidal. Then the same holds for quasi-$t$-b.l. sections.
\begin{lemma}[QuasiNF]
  \label{lem:quasi-nf}
  \lean{BSC.TODO.QuasiNF}
  \uses{def:quasi-tbirl, lem:nf-lift}
  \textit{Source: Bresciani, arXiv:2108.13397, Lemma 10.}
  Let \(k\) be a number field. If every \(\TBL\) section over \(k\) is geometric or
  cuspidal, then every quasi-\(\TBL\) section over \(k\) is geometric or cuspidal.
\end{lemma}

% SOURCE QUOTE PROOF: Let $X$ be a curve over $k$ and $s\in\Pi_{X/k}(k)$ a quasi-$t$-b.l.
% section. We may assume that there exists a non-parabolic curve $Y$ and a dominant
% morphism $f:X\to Y$ such that $f(s)$ is $t$-b.l. Since $s$ is b.l., for every finite
% place $\nu$ of $k$ there exists a unique associated point $x_{\nu}\in \bar{X}(k_{\nu})$.
% Since $f(s)$ is $t$-b.l., there exists a unique $k$-rational point $y\in \bar{Y}(k)$
% associated to $f(s)$, in particular $f(x_{\nu})=y$ for every $\nu$. Any b.l. lifting
% of $s$ to $X\setminus f^{-1}(y)$ is cuspidal thanks to \cite[Theorem B]{sti15}, hence
% $s$ is geometric or cuspidal.
\begin{proof}
  \uses{lem:nf-lift}
  Let \(s\) be quasi-\(\TBL\) with dominant \(f : X \to Y\) (\(Y\) non-parabolic) and
  \(f(s)\) \(\TBL\). Since \(f(s)\) is \(\TBL\), there is a unique \(k\)-rational
  point \(y \in \bar{Y}(k)\) with \(f(x_\nu) = y\) for every finite place \(\nu\).
  Any \(\BL\) lifting of \(s\) to \(X \setminus f^{-1}(y)\) is cuspidal by Stix's
  Theorem~B; hence \(s\) is geometric or cuspidal.
\end{proof}

% SOURCE: [Bresciani, arXiv:2108.13397], Proposition 3 (Tamagawa), p.~18
% SOURCE QUOTE: Let $X$ be a hyperbolic curve over a field $k$ finitely generated over
% $\Q$ and let $s\in\Pi_{X/k}(k)$ be a Galois section. If $s$ is not geometric nor
% cuspidal, there exists a curve $Y$ of genus $\ge 2$ with $\bar{Y}(k)=\emptyset$, a
% finite étale morphism $Y\to X$ and a lifting $r\in\Pi_{Y/k}(k)$ of $s$.
\begin{proposition}[TamagawaLem]
  \label{prop:tamagawa}
  \lean{BSC.TODO.TamagawaLem}
  \textit{Source: Bresciani, arXiv:2108.13397, Proposition 3 (Tamagawa).}
  Let \(X\) be a hyperbolic curve over \(k\) finitely generated over \(\Q\) and
  \(s \in \FundGerbe{X}{k}(k)\) a Galois section. If \(s\) is neither geometric nor
  cuspidal, there exists a curve \(Y\) of genus \(\geq 2\) with \(\bar{Y}(k) = \emptyset\),
  a finite étale morphism \(Y \to X\), and a lifting \(r \in \FundGerbe{Y}{k}(k)\) of
  \(s\).
\end{proposition}

\begin{proof}
  This is Tamagawa (1997), Proposition 2.8(iv); see also Stix's Lemma 53 in
  \textit{Rational Points and Arithmetic of Fundamental Groups} (2013).
\end{proof}

% SOURCE: [Bresciani, arXiv:2108.13397], Proposition 4 (nffg), p.~19
% SOURCE QUOTE: Assume that $t$-b.l. sections of curves defined over number fields are
% geometric or cuspidal. Then quasi-$t$-b.l. sections of curves over fields of finite
% type over $\Q$ are geometric or cuspidal.
\begin{proposition}[NFFG]
  \label{prop:nffg}
  \lean{BSC.TODO.NFFG}
  \uses{lem:quasi-nf, prop:tamagawa, lem:spec-quasi, lem:morph-birl}
  \textit{Source: Bresciani, arXiv:2108.13397, Proposition 4.}
  Assume that \(\TBL\) sections of curves over number fields are geometric or cuspidal.
  Then quasi-\(\TBL\) sections of curves over fields of finite type over \(\Q\) are
  geometric or cuspidal.
\end{proposition}

% SOURCE QUOTE PROOF: We prove this by induction on the transcendence degree $n$ of the
% base field $k$ over $\Q$. The case $n=0$ is \autoref{quasinf}. Assume $n \ge 1$ and
% that the statement is proved for fields of transcendence degree $n-1$, let $h\s k$ be
% algebraically closed in $k$ and of transcendence degree $n-1$ over $\Q$. [...] assume
% by contradiction that $s$ is not geometric nor cuspidal, thanks to Tamagawa's argument
% (\autoref{tama}) and \autoref{morb} up to passing to an étale neighbourhood we may
% assume that $\bar{X}(k)=\emptyset$ [...] Let $C$ be an affine curve over $h$ whose
% function field is $k$ [...] The generic section $s$ extends to a global section
% $\tilde{s}:C\to\Pi_{\tilde{X}/C}$ thanks to \cite[Corollary 3.4]{scfg}. The
% specializations of $\tilde{s}$ are geometric thanks to \autoref{spequasi} plus
% inductive hypothesis. It follows that $s$ is geometric thanks to
% \cite[Definition 4.1, Corollary 4.9]{scfg}.
\begin{proof}
  \uses{lem:quasi-nf, prop:tamagawa, lem:spec-quasi, lem:morph-birl}
  Induct on \(\mathrm{trdeg}(k/\Q)\). The base case \(n = 0\) is \cref{lem:quasi-nf}.
  For \(n \geq 1\): assume \(s\) is not geometric/cuspidal; by \cref{prop:tamagawa}
  and \cref{lem:morph-birl} pass to \(Y\) with \(\bar{Y}(k) = \emptyset\); choose an
  elliptic curve \(E/h\) (\(h\) the algebraic closure of \(\Q\) in \(k\), trdeg
  \(n-1\)) and a finite morphism \(Y \to E_k\); take \(C/h\) affine with \(k(C) = k\)
  and spread \(s\) to \(\tilde{s} : C \to \FundGerbe{\tilde{X}}{C}\) (Bresciani,
  arXiv:2003.04649, Cor.\ 3.4); specializations are geometric by \cref{lem:spec-quasi}
  and the inductive hypothesis; conclude \(s\) is geometric by \textit{ibid.}\ Cor.\ 4.9.
\end{proof}

\section{Proof of the main theorems}

%%==========================================================================
%% PART 8: Proof of the main theorems
%%==========================================================================

% SOURCE: [Bresciani, arXiv:2108.13397], Lemma 11 (cuspbl), p.~20
% SOURCE QUOTE: Let $X$ be a curve over a countable field $k$. The following are
% equivalent.
% \begin{itemize}
%   \item For every pair of non-empty open subset $V\subset U\subset X$, the map
%     $\mc{S}_{V/k}\to\mc{S}_{U/k}$ is surjective.
%   \item For every non-empty open subset $U\subset X$, every Galois section of $U$
%     is birationally liftable.
% \end{itemize}
\begin{lemma}[CuspBL]
  \label{lem:cusp-bl}
  \lean{BSC.TODO.CuspBL}
  \uses{def:galois-section, def:birl}
  \textit{Source: Bresciani, arXiv:2108.13397, Lemma 11.}
  Let \(X\) be a curve over a countable field \(k\). The following are equivalent:
  \begin{enumerate}
    \item For every pair of non-empty open subsets \(V \subset U \subset X\), the map
      \(\GalSec{V}{k} \to \GalSec{U}{k}\) is surjective.
    \item For every non-empty open subset \(U \subset X\), every Galois section of
      \(U\) is birationally liftable.
  \end{enumerate}
\end{lemma}

% SOURCE QUOTE PROOF: The second condition clearly implies the first. Assume that the
% first holds and let $X'\subset X$ be a non-empty open subset. Since $k$ is countable,
% then $X'$ has a countable number of open subsets. Since
% $\mc{S}_{k(X')/k}\simeq\projlim_{U\subset X'}\mc{S}_{U/k}$ \cite[Lemma 259]{sti13},
% the projective system $(\mc{S}_{U/k})_{U}$ is countable and the transition maps are
% surjective, then $\mc{S}_{k(X')/k}\to\mc{S}_{X'/k}$ is surjective as well, i.e.
% every Galois section of $X'$ is birationally liftable.
\begin{proof}
  (2) \(\Rightarrow\) (1) is clear. For (1) \(\Rightarrow\) (2): since \(k\) is
  countable, \(X'\) has countably many open subsets; since
  \(\GalSec{k(X')}{k} \cong \varprojlim_U \GalSec{U}{k}\) (Stix, Lemma 259) and the
  transition maps are surjective, the limit map \(\GalSec{k(X')}{k} \to \GalSec{X'}{k}\)
  is surjective.
\end{proof}

% SOURCE: [Bresciani, arXiv:2108.13397], Corollary 4 (cuspbl2), p.~20
% SOURCE QUOTE: The following are equivalent.
% \begin{itemize}
%   \item The cuspidalization conjecture holds.
%   \item For every smooth, geometrically connected hyperbolic curve $X$ over a field
%     $k$ finitely generated over $\Q$, every Galois section of $X$ is birationally
%     liftable.
% \end{itemize}
\begin{corollary}[CuspBL2]
  \label{cor:cusp-bl2}
  \lean{BSC.TODO.CuspBL2}
  \uses{lem:cusp-bl, def:cuspidalization}
  \textit{Source: Bresciani, arXiv:2108.13397, Corollary 4.}
  The Cuspidalization Conjecture is equivalent to: for every smooth, geometrically
  connected hyperbolic curve \(X\) over \(k\) finitely generated over \(\Q\), every
  Galois section of \(X\) is birationally liftable.
\end{corollary}

% SOURCE QUOTE PROOF: Follows from \autoref{cuspbl} and the definition of the
% Cuspidalization Conjecture.
\begin{proof}
  \uses{lem:cusp-bl, def:cuspidalization}
  The Cuspidalization Conjecture is exactly Condition~(1) of \cref{lem:cusp-bl}:
  for every pair \(V \subset U \subset X\), the map \(\GalSec{V}{k} \to \GalSec{U}{k}\)
  is surjective. By \cref{lem:cusp-bl}, this is equivalent to Condition~(2): every
  Galois section of every non-empty open subset \(U \subset X\) is birationally
  liftable. Taking \(U = X\) gives the equivalence stated.
\end{proof}

% SOURCE: [Bresciani, arXiv:2108.13397], Lemma 12 (hilbsp), p.~20
% SOURCE QUOTE: Let $k$ be a Hilbertian field, $V$ an open subset of $\A^{1}$, $X\to V$
% a family of curves, $s_{1},s_{2}\in\Pi_{X/V}(V)$ two sections. If $s_{1},s_{2}$ have
% isomorphic specializations at $k$-rational points of $V$, then $s_{1}\simeq s_{2}$.
\begin{lemma}[HilbertSections]
  \label{lem:hilbert-sections}
  \lean{BSC.TODO.HilbertSections}
  \uses{lem:rel-gerbe-limit}
  \textit{Source: Bresciani, arXiv:2108.13397, Lemma 12.}
  Let \(k\) be a Hilbertian field, \(V\) an open subset of \(\A^1\),
  \(X \to V\) a family of curves, and \(s_1, s_2 \in \FundGerbe{X}{V}(V)\) two
  sections. If \(s_1\) and \(s_2\) have isomorphic specializations at all \(k\)-rational
  points of \(V\), then \(s_1 \cong s_2\).
\end{lemma}

% SOURCE QUOTE PROOF: Write $\Pi_{X/V}=\projlim\Phi_{i}$ as in \autoref{limit}, let
% $W_{i}$ be the fibered product $V\times_{\Phi_{i}}V$ with respect to the two
% morphisms $s_{1},s_{2}:V\to\Pi_{X/V}\to\Phi_{i}$, we have that
% $\uisom(s_{1},s_{2})=V\times_{\Pi_{X/V}}V=\projlim_{i}W_{i}$. [...] it is a finite
% étale cover since $\Phi_{i}$ is a proper étale gerbe over $V$. By hypothesis,
% $\uisom(s_{1},s_{2})(k)\to V(k)$ is surjective, this implies that
% $W_{i}(k)\to V(k)$ is surjective too. Since $k$ is Hilbertian and $W_{i}\to V$ is
% finite étale, there is a section $V\to W_{i}$. [...] $\uisom(s_{1},s_{2})(V)=
% \projlim_{i}W_{i}(V)$ is non-empty since a projective limit of finite, non-empty
% sets is non-empty.
\begin{proof}
  \uses{lem:rel-gerbe-limit}
  Write \(\FundGerbe{X}{V} = \varprojlim \Phi_i\) by \cref{lem:rel-gerbe-limit}.
  Set \(W_i = V \times_{\Phi_i} V\) (the isom-sheaf); each \(W_i \to V\) is finite
  étale. The hypothesis forces \(W_i(k) \to V(k)\) surjective; since \(k\) is
  Hilbertian, each \(W_i(V) \neq \emptyset\). A projective limit of finite nonempty
  sets is nonempty, so \(\underline{\mathrm{Isom}}(s_1, s_2)(V) \neq \emptyset\).
\end{proof}

% SOURCE: [Bresciani, arXiv:2108.13397], Lemma 13 (bireq), p.~22
% SOURCE QUOTE: The following are equivalent.
% \begin{itemize}
%   \item The birational section conjecture holds.
%   \item For every curve $X$ over a field $k$ finitely generated over $\Q$, every
%     b.l. Galois section of $X$ is either geometric or cuspidal.
% \end{itemize}
\begin{lemma}[BiratEquiv]
  \label{lem:bireq}
  \lean{BSC.TODO.BiratEquiv}
  \uses{def:birl, def:bsc}
  \textit{Source: Bresciani, arXiv:2108.13397, Lemma 13.}
  The following are equivalent:
  \begin{enumerate}
    \item The BSC holds.
    \item For every curve \(X\) over \(k\) finitely generated over \(\Q\), every
      \(\BL\) Galois section of \(X\) is geometric or cuspidal.
  \end{enumerate}
\end{lemma}

% SOURCE QUOTE PROOF: The first condition clearly implies the second. Assume that the
% second holds [...] For every open subset $U\subset X$, denote by $\iota_{U}(s)$ the
% image of $s$ in $\mc{S}_{U/k}$, by hypothesis it is cuspidal. By
% \cite[Theorem 17 (2)]{sti12}, there exists a unique rational point $x\in\bar{X}(k)$
% associated with $\iota_{X}(s)$ [...] For every open subset $U\subset X$, denote by
% $\mc{P}_{U,x}\subset\mc{S}_{U/k}$ the packet of cuspidal sections associated with
% $x$ [...] $s$ is the limit of the sections $\iota_{U}(s)$, we get that $s$ is
% cuspidal associated with $x$ too.
\begin{proof}
  (1) \(\Rightarrow\) (2) is clear. For (2) \(\Rightarrow\) (1): let
  \(s \in \GalSec{k(X)}{k}\) be birational; for every open \(U \subset X\),
  \(\iota_U(s)\) is \(\BL\), hence cuspidal by hypothesis. By Stix, Theorem 17(2),
  there is a unique \(x \in \bar{X}(k)\) associated with \(\iota_X(s)\). Since
  \(\GalSec{k(X)}{k} \cong \varprojlim_U \GalSec{U}{k}\) and \(\iota_U(s) \in \mathcal{P}_{U,x}\)
  for every \(U\), we get \(s \in \varprojlim_U \mathcal{P}_{U,x}\), so \(s\) is
  cuspidal.
\end{proof}

% SOURCE: [Bresciani, arXiv:2108.13397], Theorem A, p.~3
% SOURCE QUOTE: Let $X$ be a smooth curve over a field $k$ finitely generated over $\Q$.
% A Galois section of $X$ is geometric or cuspidal if and only if it is $t$-birationally
% liftable.
\begin{theorem}[TheoremA]
  \label{thm:main-A}
  \lean{BSC.TODO.TheoremA}
  \uses{def:tbirl, def:geometric-section, def:cuspidal-section, def:galois-section}
  \textit{Source: Bresciani, arXiv:2108.13397, Theorem A.}
  Let \(X\) be a smooth curve over a field \(k\) finitely generated over \(\Q\). A
  Galois section of \(X\) is geometric or cuspidal if and only if it is \(\TBL\).
\end{theorem}

% SOURCE QUOTE PROOF: Thanks to \autoref{nffg}, we may assume that $k$ is a number field.
% Let $s\in\Pi_{X/k}(k)$ be a $t$-b.l. section, assume by contradiction that $s$ is
% neither geometric nor cuspidal. By \autoref{nflift} we may assume that $X$ is
% hyperbolic. By Tamagawa's argument (\autoref{tama}) and \autoref{mortb}, up to passing
% to an étale neighbourhood of $s$ we may assume that $X$ has genus at least $2$ and
% $\bar{X}(k)=\emptyset$. Fix $\nu$ any finite place of $k$ and let
% $x_{\nu}\in \bar{X}(k_{\nu})$ the local point associated with $s$. Choose a projective
% embedding $j:\bar{X}\s\P^{n}$ such that $j(x_{\nu})\in\A^{n}\s\P^{n}$ and let
% $U=\A^{n}\cap X$ [...] $c(j(s'))\in\Pi_{\A^{1}_{k}\setminus\{1,2\}/k}(k)$ is a
% $t$-b.l. section whose base change to $k_{\nu}$ is associated with a non-$k$-rational
% local point, which is in contradiction with \autoref{simple}.
\begin{proof}
  \uses{prop:nffg, lem:nf-lift, prop:tamagawa, cor:morph-tbirl, prop:simple}
  By \cref{prop:nffg} reduce to \(k\) a number field. Suppose \(s\) is \(\TBL\) but
  not geometric/cuspidal. By \cref{lem:nf-lift}, assume \(X\) hyperbolic. By
  \cref{prop:tamagawa} and \cref{cor:morph-tbirl}, pass to an étale neighborhood with
  \(\bar{X}(k) = \emptyset\) and \(\mathrm{genus}(X) \geq 2\). Fix a finite place
  \(\nu\) and local point \(x_\nu \in \bar{X}(k_\nu)\). Embed \(\bar{X} \subset \PP^n\)
  with \(x_\nu \in \A^n\), let \(U = \A^n \cap X\), and lift \(s\) to a \(\TBL\)
  section \(s'\) of \(U\) via \cref{lem:nf-lift}. Since \(\bar{X}(k) = \emptyset\),
  some coordinate \(c : \A^n \to \A^1\) has \(c(x_\nu)\) non-\(k\)-rational. Shrink
  \(U\) so \(c(U) \subseteq \A^1 \setminus \{1, 2\}\); then \(c(s')\) is a \(\TBL\)
  section of \(\A^1 \setminus \{1, 2\}\) whose local point is non-\(k\)-rational,
  contradicting \cref{prop:simple}.
\end{proof}

% SOURCE: [Bresciani, arXiv:2108.13397], Theorem B, p.~2
% SOURCE QUOTE: Let $X$ be a hyperbolic curve over a field $k$ finitely generated over
% $\Q$. The following are equivalent.
% \begin{itemize}
%   \item For every finitely generated extension $K/k$, every Galois section of $X_{K}$
%     is either geometric or cuspidal.
%   \item For every finitely generated extension $K/k$, every Galois section of $X_{K}$
%     is birationally liftable.
% \end{itemize}
% As a consequence, the section conjecture is equivalent to the cuspidalization
% conjecture.
\begin{theorem}[TheoremB]
  \label{thm:main-B}
  \lean{BSC.TODO.TheoremB}
  \uses{thm:main-A, cor:cusp-bl2}
  \textit{Source: Bresciani, arXiv:2108.13397, Theorem B.}
  Let \(X\) be a hyperbolic curve over \(k\) finitely generated over \(\Q\). The
  following are equivalent:
  \begin{enumerate}
    \item For every finitely generated extension \(K/k\), every Galois section of
      \(X_K\) is geometric or cuspidal.
    \item For every finitely generated extension \(K/k\), every Galois section of
      \(X_K\) is birationally liftable.
  \end{enumerate}
  As a consequence, the Section Conjecture is equivalent to the Cuspidalization
  Conjecture.
\end{theorem}

% SOURCE QUOTE PROOF: Follows directly from \hyperlink{TA}{Theorem A} and
% \autoref{cuspbl2}.
\begin{proof}
  \uses{thm:main-A, cor:cusp-bl2}
  Follows directly from \cref{thm:main-A} and \cref{cor:cusp-bl2}.
\end{proof}

% SOURCE: [Bresciani, arXiv:2108.13397], Theorem C, p.~4
% SOURCE QUOTE: The following are equivalent.
% \begin{itemize}
%   \item The birational section conjecture holds.
%   \item For every number field $k$, every section $z\in\mc{S}_{k(\P_{1})/k}$ and
%     every open subset $U\s\P^{1}$, there exists an open subset $V\s U$ and a $\pi_{1}$
%     section $r$ of $U\times V\setminus\Delta \to V$ such that the specialization
%     $r_{v}$ is equal to $\iota_{U\setminus\{v\}}(z)$ for every $v\in V(k)$.
% \end{itemize}
\begin{theorem}[TheoremC]
  \label{thm:main-C}
  \lean{BSC.TODO.TheoremC}
  \uses{lem:bireq, lem:hilbert-sections, prop:simple, lem:unique-loop-curve, lem:spec-birl, cor:cusp-bl2, prop:prorootval}
  \textit{Source: Bresciani, arXiv:2108.13397, Theorem C.}
  The following are equivalent:
  \begin{enumerate}
    \item The BSC holds.
    \item For every number field \(k\), every \(z \in \GalSec{k(\PP^1)}{k}\), and
      every open \(U \subseteq \PP^1\), there exists an open \(V \subseteq U\) and a
      \(\pi_1\)-section \(r\) of \(U \times V \setminus \Delta \to V\) such that the
      specialization \(r_v = \iota_{U \setminus \{v\}}(z)\) for every \(v \in V(k)\).
  \end{enumerate}
\end{theorem}

% SOURCE QUOTE PROOF: Thanks to \autoref{bireq}, we may use the alternative formulation
% of the birational section conjecture. Assume that the birational section conjecture
% holds and let $z$, $U$ be as in the statement. [...] By hypothesis, $z$ is associated
% with a unique rational point $p\in\P^{1}(k)$ [...] Assume that the second condition
% holds. Thanks to \cite[Theorem C]{saty21}, it is enough to consider number fields.
% [...] Passing to the limit along open subsets of $\P^{1}$ we get a Galois section of
% $\spec k(\A^{1})\otimes_{k}k(t)\setminus\{\delta\}$ which lifts $s_{k(t)}$, hence
% $x_{\nu}$ is $k$-rational by \autoref{simple}.
\begin{proof}
  \uses{lem:bireq, lem:hilbert-sections, prop:simple, lem:unique-loop-curve, lem:spec-birl, prop:prorootval}
  By \cref{lem:bireq}, use the alternative formulation. \textbf{(1) \(\Rightarrow\) (2):}
  Assume the BSC and let \(z, U\) be given. Then \(z\) is cuspidal at a unique
  \(p \in \PP^1(k)\); if \(p \in U\) take \(V = U \setminus \{p\}\) and define \(r\)
  via the diagonal. If \(p \notin U\), note \(\iota_U(z)\) is cuspidal, lift to a
  cuspidal \(z'\) of \(U_{k(t)} \setminus \{\delta\}\); all specializing loops at
  rational \(v\) are constant by \cref{lem:unique-loop-curve}, so \(z'\) extends to
  the desired \(\pi_1\)-section by \cref{prop:prorootval} and \(\cref{lem:hilbert-sections}\).
  \textbf{(2) \(\Rightarrow\) (1):} By Saidí-Tyler (Theorem~C), reduce to number
  fields; for \(b.l.\)\ \(s\) with lifting \(z\) and place \(\nu\), pass to the limit
  over open subsets using \cref{lem:hilbert-sections} to get a section of
  \(\spec k(\A^1) \otimes_k k(t) \setminus \{\delta\}\) lifting \(s_{k(t)}\); then
  \(x_\nu \in k\) by \cref{prop:simple}.
\end{proof}
